\subsection*{Forslag til opdeling af Stor Bane i to løb}
Hvis Stor Bane opdeles i to løb med de hurtigste både i løb 1 og de langsomste i løb 2, laves opdelingen automatisk fra de seneste færdighedsestimater.
Kun både med observationer i 2026 er medtaget i denne tabel.
Her bruges seneste estimate pr. båd (posteriori ved tid $t$, ellers a priori). Splitpunktet vælges, så forventet antal startende pr. race bliver så ens som muligt baseret på historisk ikke-DNC-rate.
Der er 13 både i Stor Bane, og bedste balance fås med 6 både i Stor bane 1 og 7 både i Stor bane 2 (forventet fremmøde ca. 5.67 mod 5.33 både pr. race):

\begin{table}[H]
\centering
\begin{tabular}{ll}
\toprule
Stor bane 1 (hurtigste, 6 både) & Stor bane 2 (langsomste, 7 både) \\
\midrule
Christian Grejs (\(-8.91\%\)) & Rikke M. Schnuchel (\(-0.97\%\)) \\
Axel Thomsen (\(-6.65\%\)) & Morten Juul Christensen (\(+0.03\%\)) \\
Bo Boddum (\(-3.63\%\)) & Blue X (\(+2.83\%\)) \\
Casper Lyhne (\(-3.36\%\)) & Henrik Züricho (\(+3.64\%\)) \\
Pia Olesen (\(-1.57\%\)) & Jan Ilsø Sørensen (\(+3.76\%\)) \\
Thomas Pedersen (\(-1.29\%\)) & Thomas Yde (\(+4.80\%\)) \\
 & Christian Lindholst (\(+9.42\%\)) \\
\bottomrule
\end{tabular}
\end{table}